\section{Justificación}
La presente investigación se fundamenta en la necesidad crítica de resolver la opacidad operativa de los 
modelos de aprendizaje profundo, abordando el problema desde una perspectiva interdisciplinaria que converge 
entre la Inteligencia Artificial Explicable y la computación gráfica avanzada. La relevancia de este trabajo 
se sustenta en tres dimensiones principales:

\begin{itemize}
    \item \textbf{Dimensión social y ética:} El despliegue masivo de sistemas de visión 
    artificial en entornos que impactan directamente la vida cotidiana y laboral, tales como la conducción 
    autónoma o el diagnóstico médico, exige mecanismos rigurosos de auditoría y rendición de cuentas. 
    Al abrir la «caja negra» de los algoritmos mediante representaciones volumétricas comprensibles, esta 
    investigación proporciona una base metodológica para que desarrolladores, reguladores y usuarios finales 
    puedan desmitificar el comportamiento interno de las redes, evaluar sus límites y garantizar un uso más 
    seguro y confiable de la tecnología.
    
    \item \textbf{Dimensión técnica y computacional:} Las aproximaciones convencionales para auditar capas 
    convolucionales suelen recurrir a la abstracción de datos en formatos planos bidimensionales. Este 
    enfoque obliga a los desarrolladores a generar y gestionar cientos de archivos de imagen independientes 
    (como formatos PNG individuales por cada canal) si se desea preservar la información numérica, lo que 
    fragmenta el análisis y resulta en un flujo de trabajo computacionalmente ineficiente y desarticulado. 
    Esta investigación resuelve dicha problemática al proponer un pipeline integrado que unifica la 
    información en un único entorno espacial interactivo. Para ello, se emplea el formato OpenVDB como 
    una solución técnica de conveniencia que encapsula la densidad del tensor en una estructura volumétrica 
    unificada, evitando la saturación del sistema por almacenamiento masivo de archivos planos. Esto permite 
    rescatar la continuidad espacial y la topología tridimensional nativa de los filtros convolucionales, 
    garantizando una exploración continua del dato científico sin los cuellos de botella logísticos de 
    los métodos tradicionales.
    
    \item \textbf{Dimensión metodológica e informática:} Academicamente, esta tesis construye un puente 
    directo entre dos ecosistemas de software que tradicionalmente han operado de forma aislada: las 
    herramientas de desarrollo de IA en Python y los entornos de renderizado y modelado procedural 
    tridimensional en Blender. El diseño de este pipeline procedural e interactivo dota a la comunidad de 
    ingeniería de un nuevo entorno de diagnóstico visual de vanguardia, estableciendo un precedente sobre 
    cómo las suites de gráficos 3D contemporáneas pueden transformarse en potentes plataformas de análisis, 
    pedagogía y exploración científica de datos complejos.
\end{itemize}